\scanpage{86}
\mspage{42}
de systèmes locaux, qui permet de construire une extension quotient de $(S\Pi M_G,\Gamma')$, extension du groupoïde $(\Pi M_G,\Gamma')$ par $\mathbf Z^C$, qu'on va noter
\[
\tag{117}
(S_C\Pi M_G,\Gamma').
\]
\deleted{Dans le cas où $\Gamma'=\Gamma^!_{G,G'}$, avec $G'=G/C$, on retrouve l'extension $\Pi M^*_{G,G'}$ de $\Pi M_{G,G'}$ par $\mathbf Z^C$ (tordu).
\[
\tag{118}
(S_C M_G,\Gamma_{G,G/C})\simeq\Pi M^*_{G,G/C}.
\]}

On va plutôt expliciter les données en termes de
\[
\diagram{\begin{tikzcd}[ampersand replacement=\&]
S_A\Pi M_G \arrow[r] \& (\Pi M_G,\unsure{\Gamma_G^!|\Gamma_G})
\end{tikzcd}}
\]
Si $G'\inserted{\simeq G/C}$ est une \deleted{spécialisation}\inserted{générisation} de $G$, alors on a
\[
\tag{11\unsure{9|8}}
\diagram{\begin{tikzcd}[ampersand replacement=\&]
S_A\unsure{\Pi}M_G \arrow[r] \& S_{A'}\Pi M_{G'},
\end{tikzcd}}
\]
qui passe aux quotients en
\[
\tag{120}
\diagram{\begin{tikzcd}[ampersand replacement=\&, column sep=large,row sep=large]
S_{\inserted{A}}\deleted{A}\Pi M_G \arrow[r,dash,"\gamma_{G,G'}"] \arrow[d] \& S_{A'}\Pi M_{G'} \arrow[d] \\
S_C\Pi M_G\deleted{\Gamma'} \arrow[r] \& \Pi M_{G'}
\end{tikzcd}}
\]
compatible avec l'action de $\Gamma_{G,G'}$, \unsure{donc|dont} \deleted{\unsure{que}} la deuxième flèche horizontale à son tour
\margin{85}
